\section{Why the batching ceiling is likely architectural}
\label{app:batching_architectural} The diffusion batching ceiling ($\leq 1.7\times$) is measured on two
engine-bounded stacks (Appendix~\ref{app:batching}). We argue here that
the ceiling is likely architectural, not an engine bug, by reasoning
about the mechanism. \paragraph{The KV-cache reuse argument.}
Autoregressive batching is cheap because of KV-cache reuse: when two
requests share a prompt prefix, the prefill is computed once and the KV
cache is shared. Even without prefix sharing, each request's KV cache
persists across decode steps, so the GPU alternates between requests at
low marginal cost per token. Block-diffusion invalidates both forms of reuse: \begin{enumerate}[nosep]
\item \textbf{Cross-request prefix sharing} requires that both requests
are at the same denoising step. Even if vLLM's mixed-phase crash were
fixed, requests at different $K$ steps cannot share a KV cache because
the attention mask differs per step.
\item \textbf{Within-request KV persistence} across denoising steps
fails because the denoising tokens' KV changes at every step (the mask
changes), so the per-step computation cannot reuse the previous step's
activations for the denoising positions.
\end{enumerate} The net effect: each denoising step is a fresh forward pass over the full
sequence (prompt + denoising positions), with a different attention mask.
The GPU does the same FLOPs whether batching or not, because the
attention pattern is request-specific and step-specific. This is
structurally different from AR decoding, where the attention pattern is
causal and the KV cache grows monotonically. Disabling prefix caching on
the AR arm drops its multiplier from $\cacheOnMult\times$ to
$\cacheOffMult\times$ --- still more than twenty times the diffusion
ceiling. A patched vLLM might raise the DG multiplier to $\sim 2$--$3\times$,
still an order of magnitude below AR. This is a structural argument, not
a measurement.
